\bookchapter{Kapitel 7: Spitzentechnologien im Detail}

\chapterintro{Die Technik entwickelt sich täglich weiter, und auch die Smartphone-Funktionen evolveieren ständig. Satellitenkommunikation, Off-Grid-Kommunikation, Unfallerkennung, SparkLink-Technik … diese nach Science-Fiction klingenden Funktionen sind bereits in unser Leben getreten. Lass uns gemeinsam diese fortschrittlichen Technologien kennenlernen.}

\booksection{Eins. Satellitenkommunikation: Ohne Empfang trotzdem um Hilfe rufen}

Die Satellitenkommunikation ist ein großer Durchbruch im Smartphone-Bereich der letzten Jahre. Wenn du in der Wildnis, tiefen Wäldern bist und dein Smartphone überhaupt kein Signal hat, kann die Satellitenkommunikation die Rettungsleine sein.

\booksubsection{Funktionsprinzip:}

Das Smartphone verbindet sich direkt mit einem Kommunikationssatelliten in der Umlaufbahn und sendet Notrufnachrichten. Derzeit wird hauptsächlich das Senden von SMS und Standortinformationen unterstützt, nicht aber Telefonate.

\booksubsection{Unterstützte Modelle:}

\begin{itemize}
  \item \textbf{Apple}: iPhone 14-Serie und neuer, unterstützt Notfall-SOS per Satellit. Die iPhone 16-Serie unterstützt zusätzlich das Senden von iMessage und SMS per Satellit.
  \item \textbf{Huawei}: Mate 50-Serie unterstützt Beidou-Satelliten-Nachrichten (nur Text), Mate 60 Pro-Serie zusätzlich Tiantong-Satellitentelefonie (Telefonate möglich), Mate 70-Serie mit weiter verbesserten Satellitenkommunikationsfähigkeiten.
  \item \textbf{Xiaomi/OPPO/vivo}: Ab 2024 unterstützen schrittweise Flaggschiff-Modelle die Beidou-Satelliten-Nachrichtenfunktion.
  \item \textbf{Trend}: Die Satellitenkommunikation entwickelt sich von einer exklusiven High-End-Funktion zu einer breiteren Verfügbarkeit; künftig könnten auch Mittelklasse-Modelle unterstützt werden. Das chinesische Beidou-Satellitensystem wird am breitesten von inländischen Smartphones unterstützt.
\end{itemize}

\booksubsection{Anwendungsszenarien:}

\begin{itemize}
  \item Outdoor-Abenteuer, Bergunglücke;
  \item Seenotfälle;
  \item Naturkatastrophen mit Unterbrechung der Bodenkommunikation.
\end{itemize}

\booksubsection{Verwendung:}

\begin{enumerate}
  \item Stelle sicher, dass du dich an einem offenen Ort mit Sicht auf den Himmel befindest.
  \item : Spitzenrate von 900 Mbit/s pro Verbindung, ein Vielfaches von Bluetooth-Audio; kann verlustfreies Audio und sogar HD-Video übertragen.
  \item Niedriger Stromverbrauch
  \item : Ab der iPhone 14-Serie wird Notruf-SOS per Satellit unterstuetzt. Die iPhone 17-Serie unterstuetzt zusaetzlich iMessage und SMS per Satellit und hat eine erweiterte regionale Abdeckung.
\end{enumerate}

\begin{tcbinfo}{Ultra}
Die Satellitenkommunikation erfordert eine spezielle Hardware-Antenne im Smartphone – nicht alle Modelle unterstützen sie. Außerdem ist für die Satellitenkommunikation normalerweise ein kostenpflichtiges Abonnement erforderlich, die Notruffunktion ist jedoch in der Regel kostenlos.

\end{tcbinfo}

\booksection{Zwei. Off-Grid-Kommunikation: Neue Art der Nahfeldkommunikation}

„Off-Grid“ bezeichnet die direkte Kommunikation zwischen zwei Smartphones, ohne dass ein Mobilfunknetz oder WLAN vorhanden ist.

\booksubsection{Technisches Prinzip:}

Nutzung von Bluetooth oder speziellen Funkfrequenzen für eine Punkt-zu-Punkt-Verbindung über kurze Distanz (üblicherweise einige hundert Meter bis wenige Kilometer).

\booksubsection{Anwendungsszenarien:}

\begin{itemize}
  \item \textbf{Outdoor-Aktivitäten}: Kommunikation zwischen Teammitgliedern beim Wandern oder Campen.
  \item \textbf{Katastrophenhilfe}: Bei Kommunikationsausfällen nach Erdbeben oder Überschwemmungen.
  \item \textbf{Großveranstaltungen}: Bei schlechtem Empfang wegen vieler Menschen auf Konzerten oder Sportevents.
\end{itemize}

\booksubsection{Unterstützte Plattformen:}

\begin{itemize}
  \item \textbf{Apple}: iOS 17 „NameDrop" und „Freunde offline finden".
  \item \textbf{Huawei}: Die MeeTime-App unterstützt Off-Grid-Anrufe.
  \item \textbf{Xiaomi}: Teile der Xiaomi-MiaoXiang-Zentrums-Funktionen.
\end{itemize}

\booksection{Drei. Unfallerkennung: Lebensrettung im entscheidenden Moment}

Suche nach „Unfallerkennung“ / „Notfall-SOS” und prüfe, ob es aktiviert ist.

\booksubsection{Funktionsprinzip:}

Das Smartphone nutzt mehrere Sensoren – Beschleunigungssensor, Gyroskop, Barometer, Mikrofon – um umfassend zu beurteilen, ob ein schwerer Aufprall stattgefunden hat. Nach Erkennung eines Unfalls:

\begin{enumerate}
  \item Es erscheint eine Countdown-Oberfläche, die fragt, ob du in Ordnung bist;
  \item Wenn der Countdown endet und du nicht reagierst, wird automatisch ein Notruf getätigt;
  \item Dein Standort wird an Notfallkontakte gesendet.
\end{enumerate}

\booksubsection{Unterstützte Modelle:}

\begin{itemize}
  \item \textbf{Apple}Niedriger Stromverbrauch
  \item \textbf{Android}Ultra
\end{itemize}

\booksubsection{Hinweise:}

\begin{itemize}
  \item Xiaomi
  \item Fehlauslösungen möglich (z. B. Achterbahn, intensive Bewegung) – einfach rechtzeitig abbrechen.
  \item Abhängig von Sensoren und Algorithmen, nicht 100 \% zuverlässig, aber eine zusätzliche Sicherheitsebene.
\end{itemize}

\booksection{NearLink (Xinglian): die drahtlose Verbindung der nächsten Generation}

Vier. SparkLink: Drahtlose Verbindung der nächsten Generation

\booksubsection{Technische Merkmale:}

\begin{itemize}
  \item \textbf{Ultra-niedrige Latenz}: End-zu-End-Latenz von 20 Mikrosekunden (0,02 ms), zwei Größenordnungen unter Bluetooth – Spiel- und Musik-Synchronisation nahezu unmerklich.
  \item \textbf{Hohe Rate}: Spitzenrate von 900 Mbit/s pro Verbindung, ein Vielfaches von Bluetooth-Audio; kann verlustfreies Audio und sogar HD-Video übertragen.
  \item \textbf{Niedriger Stromverbrauch}: Standby-Stromverbrauch vergleichbar mit Bluetooth, aber höhere Übertragungseffizienz und geringerer Energieverbrauch pro Dateneinheit.
  \item \textbf{Android}: Verbindungsreichweite über 60 \% größer als Bluetooth; ein Netzwerk unterstützt über tausend gleichzeitig verbundene Geräte.
\end{itemize}

\booksubsection{Hinweise:}

\begin{itemize}
  \item \textbf{Kabellose Kopfhörer}: Geringere Latenz, perfekte Audio-Video-Synchronisation beim Spielen.
  \item \textbf{Smartwatches}: Stabilere Verbindung.
  \item \textbf{Smart Home}: Mehr Geräte gleichzeitig verbunden.
  \item \textbf{Automobil-Vernetzung}: Flüssigere Handy-Auto-Verbindung.
\end{itemize}

\booksubsection{Unterstützte Geräte:}

Aktuell hauptsächlich Huawei-Ökosystem-Geräte, z. B. Huawei Mate 60-Serie, bestimmte Huawei-Lautsprecher und -Watches. Künftig werden weitere Marken folgen.

\textbf{Technischer Ausblick}\par

SparkLink repräsentiert die Zukunft der drahtlosen Kommunikation. Mit dem Beitritt weiterer Hersteller werden wir voraussichtlich mehr SparkLink-fähige Geräte sehen und ein besseres drahtloses Erlebnis genießen.

\begin{tcbpractice}{Prüfe: Wird dein Smartphone unterstützt?}
\begin{enumerate}
  \item Suche in den Einstellungen nach „Satellit" und prüfe, ob Optionen zur Satellitenkommunikation vorhanden sind.
  \item Suche nach „Unfallerkennung" / „Notfall-SOS" und prüfe, ob es aktiviert ist.
  \item Wenn du ein Huawei-Smartphone hast, prüfe, ob SparkLink-Geräte unterstützt werden.
  \item Lerne die Verwendung dieser Funktionen kennen – sie könnten im entscheidenden Moment nützlich sein.
\end{enumerate}

\end{tcbpractice}
